% Wygenerowano z ../01_lidar.md skryptem _md_do_latex.py (można edytować ręcznie).
\section{Ćwiczenie 1: LiDAR, czyli jak robot mierzy odległość (12 min)}\label{ux107wiczenie-1-lidar-czyli-jak-robot-mierzy-odlegux142oux15bux107-12-min}

\textbf{Notatnik:} \texttt{01\_lidar.ipynb}

\textbf{Potrzebne:} miarka, karton (ok. 30 × 30 cm), zestaw materiałów: biała kartka, czarna matowa tkanina, szyba lub plexi, lusterko, metalowa puszka, kątomierz i linijka (do ustawiania kątów)

\subsection{Cel}\label{cel}

Sprawdzić, jak dokładnie LiDAR mierzy odległość i w jakich sytuacjach się myli.

\subsection{Trochę teorii}\label{trochux119-teorii}

LiDAR (\emph{Light Detection and Ranging}) wysyła krótkie impulsy światła laserowego i mierzy, po jakim czasie wracają po odbiciu od przeszkody. Ponieważ światło porusza się ze stałą prędkością, z czasu powrotu można policzyć odległość. Głowica LiDAR-u obraca się kilka razy na sekundę. W ten sposób robot dostaje \textbf{kilkaset pomiarów odległości dookoła siebie}. Taki zestaw punktów to \textbf{skan} (chmura punktów).

\begin{itemize}
\tightlist
\item
  \textbf{LiDAR 2D} (taki jak na Jetbocie) mierzy w jednej płaszczyźnie, na wysokości czujnika. Nie widzi więc przedmiotów niższych lub wyższych od tej płaszczyzny.
\item
  \textbf{LiDAR 3D} ma wiele wiązek ustawionych pod różnymi kątami i widzi otoczenie w przestrzeni. Takie czujniki stosuje się w samochodach autonomicznych.
\end{itemize}

\subsection{Kroki}\label{kroki}

\subsubsection{Krok 1: pierwszy skan (2 min)}\label{krok-1-pierwszy-skan-2-min}

Uruchom komórki \texttt{polacz()} i \texttt{rysuj\_skan()}. Widzisz otoczenie robota z góry. Robot to czerwony trójkąt, a jego przód jest skierowany w górę wykresu.

\begin{itemize}
\tightlist
\item
  Gdzie na wykresie są ściany? Gdzie przeszkody na torze?
\item
  Jedna osoba staje przed robotem, a druga ponownie uruchamia \texttt{rysuj\_skan()}. Jak na skanie wyglądają nogi?
\end{itemize}

\subsubsection{Krok 2: dokładność pomiaru (4 min)}\label{krok-2-dokux142adnoux15bux107-pomiaru-4-min}

Ustaw karton \textbf{na wprost robota}, prostopadle do niego. Odległość mierz miarką od \textbf{środka LiDAR-u} (obrotowej głowicy) do kartonu. Dla każdej odległości uruchom \texttt{seria\_pomiarow(kat=0)}. Komórka wykona 10 pomiarów i poda średnią oraz rozrzut.

\par\noindent\begin{minipage}{\linewidth}
\begin{tabular}{@{}
  >{\raggedright\arraybackslash}p{(\linewidth - 6\tabcolsep) * \real{0.2366}}
  >{\raggedright\arraybackslash}p{(\linewidth - 6\tabcolsep) * \real{0.2258}}
  >{\raggedright\arraybackslash}p{(\linewidth - 6\tabcolsep) * \real{0.2581}}
  >{\raggedright\arraybackslash}p{(\linewidth - 6\tabcolsep) * \real{0.2796}}@{}}
\toprule\noalign{}
\begin{minipage}[b]{\linewidth}\raggedright
Odległość z miarki {[}m{]}
\end{minipage} & \begin{minipage}[b]{\linewidth}\raggedright
Średnia z LiDAR-u {[}m{]}
\end{minipage} & \begin{minipage}[b]{\linewidth}\raggedright
Rozrzut (odch. std.) {[}m{]}
\end{minipage} & \begin{minipage}[b]{\linewidth}\raggedright
Błąd = LiDAR − miarka {[}cm{]}
\end{minipage} \\
\midrule\noalign{}
0,30 & & & \\
0,60 & & & \\
1,00 & & & \\
1,50 & & & \\
\bottomrule\noalign{}
\end{tabular}
\end{minipage}\par\medskip

\subsubsection{Krok 3: różne materiały (3 min)}\label{krok-3-ruxf3ux17cne-materiaux142y-3-min}

Na kartonie ustawionym w odległości \textbf{0,60 m} umieszczaj kolejno różne materiały. \textbf{Lusterko} zmierz dwa razy: ustawione na wprost robota i obrócone o ok. 45°. Przy lusterku pod kątem uruchom też \texttt{rysuj\_skan()}: gdzie jest punkt?

\par\noindent\begin{minipage}{\linewidth}
\begin{tabular}{@{}
  >{\raggedright\arraybackslash}p{(\linewidth - 6\tabcolsep) * \real{0.2222}}
  >{\raggedright\arraybackslash}p{(\linewidth - 6\tabcolsep) * \real{0.2121}}
  >{\raggedright\arraybackslash}p{(\linewidth - 6\tabcolsep) * \real{0.3030}}
  >{\raggedright\arraybackslash}p{(\linewidth - 6\tabcolsep) * \real{0.2626}}@{}}
\toprule\noalign{}
\begin{minipage}[b]{\linewidth}\raggedright
Materiał
\end{minipage} & \begin{minipage}[b]{\linewidth}\raggedright
Średnia z LiDAR-u {[}m{]}
\end{minipage} & \begin{minipage}[b]{\linewidth}\raggedright
Ile pomiarów nieudanych (z 10)?
\end{minipage} & \begin{minipage}[b]{\linewidth}\raggedright
Czy wynik jest wiarygodny?
\end{minipage} \\
\midrule\noalign{}
karton (odniesienie) & & & \\
biała kartka & & & \\
czarna matowa tkanina & & & \\
szyba / plexi & & & \\
lusterko na wprost & & & \\
lusterko pod kątem 45° & & & \\
metalowa puszka & & & \\
\bottomrule\noalign{}
\end{tabular}
\end{minipage}\par\medskip

\subsubsection{Krok 4: przeszkoda pod kątem (2 min)}\label{krok-4-przeszkoda-pod-kux105tem-2-min}

Karton w odległości \textbf{0,60 m}, obracany względem robota. Kąt ustawcie kątomierzem: linijkę przyłóżcie do przedniej ściany kartonu, kątomierz do linijki (0° = karton ustawiony prostopadle do robota).

\par\noindent\begin{minipage}{\linewidth}
\begin{tabular}{@{}
  >{\raggedright\arraybackslash}p{(\linewidth - 4\tabcolsep) * \real{0.3433}}
  >{\raggedright\arraybackslash}p{(\linewidth - 4\tabcolsep) * \real{0.3134}}
  >{\raggedright\arraybackslash}p{(\linewidth - 4\tabcolsep) * \real{0.3433}}@{}}
\toprule\noalign{}
\begin{minipage}[b]{\linewidth}\raggedright
Kąt ustawienia kartonu
\end{minipage} & \begin{minipage}[b]{\linewidth}\raggedright
Średnia z LiDAR-u {[}m{]}
\end{minipage} & \begin{minipage}[b]{\linewidth}\raggedright
Pomiary nieudane (z 10)
\end{minipage} \\
\midrule\noalign{}
0° (prostopadle) & & \\
30° & & \\
60° & & \\
prawie równolegle (80°) & & \\
\bottomrule\noalign{}
\end{tabular}
\end{minipage}\par\medskip

\subsection{Pytania do przemyślenia}\label{pytania-do-przemyux15blenia}

\begin{enumerate}
\def\labelenumi{\arabic{enumi}.}
\tightlist
\item
  Czy błąd pomiaru rośnie z odległością? O ile centymetrów może się mylić LiDAR?
\item
  Przy których materiałach LiDAR zawodzi? Dlaczego? \emph{(Podpowiedź: co dzieje się ze światłem przy szybie, czarnej tkaninie i lustrze? Dlaczego lustro na wprost i pod kątem dało inne wyniki?)}
\item
  Robot autonomiczny jedzie korytarzem z przeszklonymi drzwiami. Jakie zagrożenie widzisz?
\item
  Stolik ma blat na wysokości 70 cm i cienkie nogi. Co „zobaczy'' LiDAR 2D zamontowany 15 cm nad podłogą?
\end{enumerate}

\begin{quote}
\textbf{Jeśli coś nie działa.} „Brak poprawnych pomiarów'' oznacza, że czujnik nie dostał odbicia. To też wynik, więc wpisz go do tabeli. Jeśli żaden skan się nie pojawia, uruchom \texttt{status()} i zawołaj prowadzącego.
\end{quote}
